\section{The Biblical Witness}

The Gospels distinguish demonic affliction from disease without supplying a
modern diagnostic chart. Matthew 4:24 lists the demonized, epileptics, and
paralytics separately; Mark 1:32--34 can place the sick and demonized side by
side. Conversely, symptoms in individual narratives overlap conditions that
now receive medical explanations. Narrative classification is theological
testimony, not permission to diagnose a contemporary person from resemblance.

\subsection{Words, agents, and narrative judgment}

The Greek vocabulary does not itself settle the phenomenon behind a scene.
The Gospels use forms of the designation commonly rendered ``unclean spirit''
and forms of \emph{daimonizomai}, commonly rendered by ``possessed'' or, more
transparently, ``demonized.'' Neither expression is a clinical category, and
neither tells a modern reader how to infer an invisible cause from symptoms.
The narratives make theological judgments through plot, speech, command,
restoration, and the evangelist's framing. ``Demonized'' remains a translation
choice rather than a diagnosis.

\subsection{Mark: authority, conflict, restoration}

Mark gives the most sustained architecture. At 1:21--28 Jesus commands the
unclean spirit without apparatus; the crowd asks what this authority means.
Mark 3:22--27 interprets the works through the image of binding the strong man,
placing particular expulsions within the invasion of Satan's domain. Mark
5:1--20 joins command, contested naming, destructive power, restoration, and
mission. The man's restored condition---clothed, composed, and sent back to
his own people---is narratively more important than the spirits' attempted
control of the exchange. Their speech belongs to the conflict scene; it is not
a deposit of reliable demonological information.

Mark 9:14--29 places a suffering child, an anguished father, failed disciples,
and Jesus' command within a lesson on dependent faith and prayer. Verse 29 is
also a textual and reception-history hinge. The tracked Robinson--Pierpont
Byzantine textform includes ``and fasting'' and its inline divergence note
reports that the Nestle--Aland text omits those words. The tracked
Douay--Rheims/Challoner witness also reads ``prayer and fasting,'' at Mark 9:28
in its numbering. The repository does not yet bind a separately identified
eclectic edition or a full critical apparatus. These controls establish the
existence of the variant, not its manuscript history or relative weight; an
argument resting on fasting must therefore name its textual tradition.

\subsection{Matthew and Luke: the kingdom has arrived}

Matthew organizes the controversy around the kingdom and the Spirit of God
(Mt 12:22--32); Luke says ``the finger of God'' (Lk 11:20), echoing Exodus and
making expulsion a sign that God's reign has arrived. Luke 10:17--20 subordinates
successful ministry to the disciples' names being written in heaven. The return
of the unclean spirit in Matthew 12:43--45 and Luke 11:24--26 is a warning
against an emptied life, not a technical case-history protocol.

\subsection{Acts: mission is not a manipulable name}

Acts preserves both continuity and boundary. Philip's mission includes unclean
spirits departing (Acts 8:7); Paul commands the divining spirit ``in the name
of Jesus Christ'' (16:18). Acts 19:11--20 then contrasts apostolic mission with
the sons of Sceva's attempted use of Jesus' name as a transferable technique.
The failure is not evidence that a longer formula was needed. It marks the
difference between belonging to Christ's mission and manipulating a powerful
name.

\subsection{What the canon does and does not give}

The canonical pattern is consequently christological and ecclesial before it
is procedural. Jesus' authority, the advent of God's kingdom, prayer, mission,
conversion, and freedom from fear hold the narratives together. Scripture
supports the Church's confession that Christ defeats evil; it does not itself
specify the later ministerial reservation, diagnostic process, or Roman ritual.

The same boundary is moral as well as historical. Narrative details describe
particular persons within particular theological plots. They are not a list
by which a reader may classify another person's illness, disability, trauma,
speech, or behavior. Several Gospel summaries distinguish illnesses from
hostile spirits even while placing both beneath Christ's authority. Modern
readers should not collapse categories the evangelists kept distinct, still
less translate either one directly into a retrospective clinical diagnosis.

This section does not settle later theological questions about natural and
preternatural causation. Its narrower conclusion is that the Gospel narratives
warrant neither dismissal of the evangelists' theological claims nor a
contemporary diagnosis by analogy.
